\section{Related Work}
\label{sec:related}

\textbf{ANN-to-SNN conversion.} Rate-based conversion replaces rectifiers with integrate-and-fire units and interprets firing rates as activations. Weight/threshold balancing, data-based normalisation, reset-by-subtraction and soft reset established the basic pipeline~\cite{diehl2015,rueckauer2017,sengupta2019,han2020rmp}. Deng and Gu derived a layerwise recursion and the closed-form bias that removes the systematic rate offset~\cite{deng2021}; QCFS replaced ReLU with a quantised clip--floor--shift activation so that the expected conversion error vanishes~\cite{bu2022qcfs}. FTBC added a forward temporal bias~\cite{wu2024ftbc}, and an optimal-transport view of the same correction was developed in~\cite{bojkovic2024ot}. Dual-Phase decoding showed that reading the membrane potential avoids the $1/T$ quantisation grid~\cite{dualphase2024}. TPP introduced two-phase probabilistic neurons to exploit temporal misalignment~\cite{bojkovic2025tpp}; PASCAL used inhibitory spikes for deterministic exact counting~\cite{ramesh2025pascal}; parallel spiking conversion, differential coding, adaptive bursts, channel-wise softplus thresholds, error-compensation learning and adaptive fission target the low-latency regime~\cite{hao2025parallel,huang2025dcgs,wang2025adafire,yang2025csqcfs,liu2025ecl,jiang2025fission}. Recent preprints continue to refine residual-membrane alignment, group-wise thresholds and fine-tuning~\cite{chen2026rmp,ji2026gtdl,xia2026neuro}. All of these are developed and evaluated on convolutional or attentional vision backbones with row-stochastic operators.

\textbf{Spiking graph networks.} Spiking GCNs are overwhelmingly trained directly with surrogate gradients~\cite{neftci2019,fang2023spikingjelly}, and large-scale evaluations of spiking GNNs contain no conversion method~\cite{sgnnbench2025}. One-Spike SNN converts both CNNs and text GCNs with single-spike phase coding~\cite{hwang2025onespike}, but does not analyse conversion error for message passing or compare ANN-relative energy. To our knowledge, no published work transfers a full conversion-error formulation to graph-level visual classification and evaluates it against a complete family of conversion baselines.

\textbf{Graph-structured vision.} Vision GNNs build $k$-NN graphs over image patches~\cite{han2022vig,knyazev2019}, while superpixel graph benchmarks convert MNIST and CIFAR10 images into 40--150-node graphs using SLIC and an 8-NN Gaussian-weighted adjacency~\cite{dwivedi2023,achanta2012}. These datasets ship \emph{with} their graph structure, which makes them a direct test bed: no graph construction or freezing heuristic is needed, and the propagation operator is fixed and non-negative. Classical graph backbones and semi-supervised splits provide the architectural substrate~\cite{kipf2017,velickovic2018,hamilton2017,yang2016planetoid}.

\textbf{Energy accounting.} SNN energy studies typically count synaptic operations and apply a 45\,nm CMOS table in which a 32-bit MAC costs \MACPJ\,pJ and an AC costs \ACPJ\,pJ~\cite{horowitz2014,eshraghian2023}. Because this ratio is only \MACACRatio, an advantage of $5\times$ essentially requires the SNN's AC count to fall below the ANN's MAC count. We report two conventions throughout: a primary event-driven AC-only count and a conservative count that also charges the real-valued propagation as MACs.

Spiking computation, graph representation learning and event-driven vision rest on a broader body of work~\cite{maass1997,krizhevsky2009,lecun1998,orchard2015,li2017cifardvs,kugele2020,liu2022spikeconverter,zhou2025edge,ji2026gtdl,xia2026neuro,dualphase2024,bojkovic2024ot,yang2025csqcfs,jiang2025fission,sgnnbench2025,hwang2025onespike}. These works provide the substrate for the components combined in Section~\ref{sec:method}.
